\chapter{Lie-Algebra Links as Latent Fields}
\label{comp:links}
\label{comp:lie-links}

\section{The question and the answer}
\label{comp:links-question}

The independent interaction links can now be inferred together with the agents' states. A link is
written as the exponential of a Lie-algebra generator, the generators have a normalized prior, and
their recognition law contributes to the same collective evidence bound as the agents. This is the
answer of \MT{ch:links}. It supplies a probability model and a recognition dynamics for relations
that were previously held as structural data or moved as parameters. The interaction graph itself
remains supplied: inference changes the links on its edge copies, not which agents are connected.
Everything still takes place at the single base point $c_*$.

There are two meanings of ``dynamic'' here, and they should be kept distinct. Within an inference
episode, the mean and precision of a Gaussian generator law follow first-order natural-gradient
descent. Across learning episodes, a prior couples successive generator values; the most probable
whole history obeys a second-order difference equation with Euclidean sign. The first is relaxation
of a recognition law. The second is stationarity of a history objective, including its future
episodes. Neither makes the generators into physical space-time fields.

The chapter's most recognizable physics formula is its linearized link equation. A positive
operator combines an anchor precision, circulation penalties on graph loops, an episode-difference
term, and a matter contribution, and acts on the generator field to give a matter current. Its
Proca-type form is a structural correspondence. The proved content is an operator on a finite
agent multigraph and an episode index, with a controlled nonlinear approximation near zero
generators and weak individual currents (\MT{prop:link-linearized-equation}).

\section{The picture}
\label{comp:links-picture}

Previously a receiver compared its state with a source transported through a supplied link.
Now that link is uncertain too. A successful comparison favors generators that make the transported
source explain the receiver; a Gaussian anchor discourages arbitrarily large generators; loop
penalties discourage inconsistent compositions around declared closed walks. Recognition must
balance those three effects while retaining uncertainty in the generators. This balance is the
link field equation, rather than an extra force appended to the evidence bound.

\begin{analogy}{structural}
An edge generator resembles a lattice gauge potential, its ordered exponentials resemble link
variables, and a metric defect of their loop product resembles a plaquette action. The analogy is
limited by the coordinates: generators are expressed at one reference point, and passive vertex
frame changes leave them unchanged. They are not potentials for transport along a base manifold.
On represented isometries the metric defect equals Wilson's trace density exactly; on general
linear links the fiber metric is an additional modeling datum
(\MT{prop:link-generator-gauge}, \MT{prop:link-loop-defect}).
\end{analogy}

The prior lives on a finite-dimensional vector space, not on the group with Haar measure.
This matters for noncompact groups. A Gaussian generator law is already a probability measure;
reweighting it by the exponential of a finite nonnegative loop action preserves normalizability.
It permits $\GL(K,\R)$ without a finite Haar reference, at the price of supplied anchors,
covariances, and an inner product. Normalization establishes the model; it does not establish
that inferred links improve prediction or evidence over held links.

\section{Objects and notation}
\label{comp:links-notation}

\begin{decoder}
$E_+$, $\mathscr T=\mathfrak g^{E_+}$ & One orientation of each reversal pair; the vector space of generators on those copies. & Oriented lattice edges; edge fields.\\
$\mathring\omega^x_e$, $\Theta^x_e$ & Reference generator and the link it determines in the frames of its endpoints. & Exponential coordinates of a link variable.\\
$\mathring\Theta^x_p$ & Ordered product of reference twists around a walk. & Based Wilson-loop matrix.\\
$\mathsf G^{\mathrm{fib}}$, $\mathfrak Q^x_p$ & Supplied fiber metric and squared deviation of represented loop holonomy from identity. & Metric plaquette defect.\\
$\mathbf C_p$ & Signed sum of generators along a closed walk. & Linearized circulation.\\
$\mathring\alpha$, $\mathring\Sigma$ & Anchor and covariance of the Gaussian base prior. & Background preference and bare stiffness.\\
$R=\mathcal N(\mathring\mu,\mathring\Lambda^{-1})$ & Recognition law of generators, with inferred mean and precision. & Variational Gaussian field law.\\
$\mathbf H_o$, $\mathbf J^x_e$ & Negative log held-link joint density; minus its expected derivative with respect to a represented twist. & Matter energy and link current.\\
$\mathbf L(\mathbf B,\mathbf T)$ & Fr\'echet derivative of the matrix exponential in direction $\mathbf T$. & Differential of exponential coordinates.\\
$n$, $\mathbf m$ & Learning-episode index and increment precision in the history prior. & Euclidean lattice index and inertial coefficient.\\
$\mathbf K$, $\mathbf N_0$, $\mathbf j$ & Positive leading operator, current-dependent Hessian correction, and projected source. & Proca-type operator and source, structurally.\\
\end{decoder}

A ring accent means reference coordinates, not a derivative. Prior covariance $\mathring\Sigma$
and recognition covariance $\mathring\Lambda^{-1}$ are different objects; ordinary $\Sigma_j$
below is an agent's matter covariance. The two channels can have independent generator families.
A common-link specialization uses one family for both. The displays suppress the channel index
when one family suffices. The reversal law supplies an inverse link even if the source-label
prior supports attention in only one direction; it does not require reciprocal attention.

\section{The formulas, explained}
\label{comp:links-formulas}

\subsection{Coordinates, gauge, and the extent of the chart}

\begin{keyequation}{From generators to independent links}
\[
\begin{aligned}
\mathring\omega^x_{\bar e}&=-\mathring\omega^x_e,\\
\Theta^x_e&=U_i\exp(\mathring\omega^x_e)U_j^{-1},\qquad e:j\to i,\\
\mathring\Theta^x_p&=\prod_{a=0}^{r-1}\exp(\mathring\omega^x_{e_a}).
\end{aligned}
\]
\end{keyequation}

The endpoint frames convert a reference twist into a map from the source frame to the receiver
frame. The loop product is ordered: noncommuting generators cannot be replaced by their sum.
Under passive vertex changes $U_i\mapsto g_iU_i$, the links transform by their usual vertex law
and every generator stays fixed. A reference change $\sigma\mapsto\sigma\cdot g_0$ instead
sends all generators through $\operatorname{Ad}_{g_0^{-1}}$. These are different operations
(\MT{def:link-exponential-chart}, \MT{prop:link-generator-gauge}).

``Chart'' does not mean a global one-to-one coordinate system. A real exponential has positive
determinant, and even $\GL^+(K)$ contains nonexponential matrices: for $K\geq2$,
$\operatorname{diag}(-1,-2,1,\ldots,1)$ has no real logarithm. The model therefore restricts
each twist to the exponential image. Products of such twists are less restricted: if a closed
walk traverses two distinct reversal pairs exactly once each, its reference holonomies reach
all of $\GL^+(K)$. For $SO(K)$ each link has a generator, though not a unique one
(\MT{prop:link-chart-image}).

\subsection{The metric defect and the normalized prior}

\begin{keyequation}{A normalizable link prior, including noncompact groups}
\[
\begin{aligned}
\mathfrak Q^x_p&=\tfrac12\lVert I-\rho_x(\mathring\Theta^x_p)\rVert_{\mathsf G}^2,
\qquad\mathscr S=\sum_p\mathrm c_p\mathfrak Q^x_p,\\
\nu^{\mathscr T}&=\bigotimes_{e\in E_+}\mathcal N(\mathring\alpha^x_e,\mathring\Sigma^x_e),\\
\mathbb P^{\mathscr T}_\theta&=Z_{\mathscr S}^{-1}e^{-\mathscr S}\nu^{\mathscr T},\\
Z_{\mathscr S}&=\E_{\nu^{\mathscr T}}e^{-\mathscr S},\qquad0<Z_{\mathscr S}\leq1.
\end{aligned}
\]
\end{keyequation}

Here $\mathrm c_p\geq0$, and the norm is induced by
$\langle\mathbf H,\mathbf H'\rangle_{\mathsf G}
=\Tr((\mathsf G^{\mathrm{fib}})^{-1}\mathbf H^\top\mathsf G^{\mathrm{fib}}\mathbf H')$.
The defect vanishes exactly at identity represented holonomy; a representation with a kernel
detects triviality in that representation. Its Gaussian expectation is finite, and
$Z_{\mathscr S}\geq\exp(-\E_{\nu^{\mathscr T}}\mathscr S)$
(\MT{prop:link-loop-defect}, \MT{prop:link-generator-prior}). With no penalized loops,
the prior is the Gaussian base law. On a translation group the defect is exactly
$\tfrac12\lVert\mathbf C_p\mathring\omega\rVert_{\mathsf G}^2$. On a general group the quadratic part is
$\tfrac12\lVert\mathrm d\rho_x((\mathbf C\mathring\omega)_p)\rVert_{\mathsf G}^2$ near zero,
with an explicit cubic error bound
(\MT{prop:link-translation-instance}, \MT{prop:link-loop-defect}).

The metric is necessary for the stated measurement on $\GL(K,\R)$. Every continuous class
function has its identity value on every unipotent matrix, including a nontrivial shear.
Consequently no such function can be a defect vanishing only at identity
(\MT{prop:link-class-function-blindness}, \MT{cor:link-no-invariant-defect}).
The metric defect is passively invariant because the reference twist is. Under a reference
change the same geometric metric must be carried along. Holding its matrix fixed preserves
all defects under changes acting by multiples of metric isometries, and, in the full general
linear setting specified by the proposition, only those. A cyclic change of loop base vertex
need not preserve the defect either. This is a supplied geometric measurement, not a
conjugation-invariant scalar on the unrestricted group.

\subsection{Where the link sector enters the evidence bound}

Given generators, use the held-link latent law and normalized record kernels. Integrating
the resulting joint over generators, matter, and records gives one, even on a cyclic record
hypergraph (\MT{prop:link-latent-normalization}). For bounded record densities and positive
conditional evidence, disintegrate recognition as $Q(d\mathring\omega,dy)
=R(d\mathring\omega)Q_{\mathring\omega}(dy)$. Then:

\begin{keyequation}{The collective free energy with uncertain links}
\[
\begin{aligned}
\Fenergy_V[Q;X,o]
={}&\KL(R\Vert\mathbb P^{\mathscr T}_\theta)
+\E_R\Fenergy_V[Q_{\mathring\omega};X,\Theta(\mathring\omega),o],\\
\KL(R\Vert\mathbb P^{\mathscr T}_\theta)
={}&\KL(R\Vert\nu^{\mathscr T})+\E_R\mathscr S+\log Z_{\mathscr S}.
\end{aligned}
\]
\end{keyequation}

The second identity uses finite base divergence and integrable action; both hold for a
nondegenerate Gaussian $R$ (\MT{prop:link-gaussian-recognition-vfe}). The first term charges
recognition for changing the link law away from its prior. The second averages the ordinary
held-link free energy over that uncertainty. The loop action enters with its normalizer;
dropping $\log Z_{\mathscr S}$ changes the model, and matters when learning prior parameters.
Correlated Gaussian generators are allowed; their divergence from the edge-product base law
includes their total correlation.

\subsection{The field equation and its delta-rule current}

For product recognition $R\otimes Q^{\mathsf Y}$, hold the matter law fixed and write
$\mathbf H_o$ for the negative log density of the full held-link joint. It includes the agent
latent law as well as record energies. Under the differentiability and integrable growth
conditions of \MT{prop:link-field-equation}, Gaussian stationarity requires:

\begin{keyequation}{Prior stress balances the expected matter current}
\[
\begin{aligned}
\mathring\Sigma^{-1}(\mathring\mu-\mathring\alpha)+\E_R\nabla\mathscr S
&=-\E_{R\otimes Q^{\mathsf Y}}\nabla_{\mathring\omega}\mathbf H_o,\\
\mathring\Lambda
&=\mathring\Sigma^{-1}+\E_R\nabla^2\mathscr S
+\E_{R\otimes Q^{\mathsf Y}}\nabla^2_{\mathring\omega}\mathbf H_o.
\end{aligned}
\]
\end{keyequation}

The anchor restores the mean, loops supply expected circulation stress, and matter supplies
a current. Precision balances expected curvature, rather than being a free fixed learning
rate. For edge-product recognition only the diagonal precision blocks are varied. The
right side is a sum of record derivatives only if the agent latent law does not read the links.

If matter reads an edge generator through $\mathbf M=\rho_x(\exp\mathring\omega^x_e)$,
its current is $\mathbf J^x_e=-\E_{Q^{\mathsf Y}}\partial\mathbf H_o/\partial\mathbf M$.
The edge's mean balance has on its right side
\[
\E_R\Bigl[(\mathrm d\rho_x)^\top
\mathbf L\bigl((\mathrm d\rho_x\mathring\omega^x_e)^\top,\mathbf J^x_e\bigr)\Bigr],
\qquad
\mathbf L(\mathbf B,\mathbf T)=\int_0^1e^{t\mathbf B}\mathbf Te^{(1-t)\mathbf B}dt.
\]
The transpose and adjoint carry a matrix current back to generator coordinates, and the
expectation averages it over the generator law. Replacing that expectation by its value at
the mean is an additional approximation, not the displayed stationarity equation.

For a binary success comparison with energy
$\tfrac12(k_i-\mathbf Mk_j)^\top W_a(k_i-\mathbf Mk_j)$, the exact current is
\[
\mathbf J^b_e=W_a\bigl[(\mu_i-\mathbf M\mu_j)\mu_j^\top
+\Sigma_{ij}-\mathbf M\Sigma_j\bigr].
\]
The first term is prediction error times source mean, the delta-rule form; the remaining
terms retain cross covariance and source uncertainty (\MT{cor:link-delta-rule}). For product
matter $\Sigma_{ij}=0$. Selective comparison instead includes the source weight and query/key
maps of \MTeq{eq:link-delta-current-attention}, with its stated generator-independent
recognition row. The correspondence with fast-weight update rules is structural, not an
identification of architectures.

Without a generator prior, product matter drives this comparison's matrix optimum to
$\mathbf M^\star=\mu_i\mu_j^\top(\Sigma_j+\mu_j\mu_j^\top)^{-1}$, of rank at most one.
For $K\geq2$ it is singular: the infimum over $\GL^+(K)$ is unattained, and the group-level
energy has no stationary point. This does not mean that generator coordinates have no
stationary points. Critical points of the exponential map can create stationary generators
whose links are not stationary for the matrix energy. Correlated matter can also give an
invertible regression optimum. At the law level, stabilizers create a different degeneracy,
nonisolated stationary links (\MT{prop:link-law-level-energy},
\MT{cor:link-prior-necessary}). With the anchor prior and lower-bounded matter objective,
the most-probable-link and Gaussian recognition problems attain minima. This proves
existence under the stated conditions, not global uniqueness.

\subsection{Held limits and induced actions}

Shrinking the prior covariance as $\mathring\Sigma=\mathring\sigma^2\mathring\Sigma_0$
recovers held anchor links in free-energy value:
\[
\inf\Fenergy_{\mathring\sigma}\longrightarrow
\inf_{Q\in\mathcal Q'}\Fenergy_V[Q;X,\Theta(\mathring\alpha),o].
\]
This assumes bounded record densities, positive conditional evidence at every generator,
a finite lower-semicontinuous held-link infimum at the anchor, and a continuous,
growth-controlled approximating subclass (\MT{thm:link-held-limit}). Bounded-energy generator
recognition concentrates there; near-minimizers also have divergence from the narrowing prior
tending to zero. Conditional evidence continuity gives convergence of evidence and posterior
concentration. Narrowing recognition alone against a fixed-width prior instead sends its
divergence to infinity; a point mass is never a finite-energy recognition optimum.

Held-link statements remain valid conditionally at every admitted generator. They transfer
to link-marginal laws only when the property survives mixtures: conditional Gaussian matter
can have a non-Gaussian marginal (\MT{prop:link-conditional-transfer}). Integrating out matter
also depends on how the joint was built. Normalized record extension gives no induced prior
action. Joint normalization of undirected potentials gives
$\mathscr S^{\mathrm{ind}}(\mathring\omega)=-\log Z_{X,\Theta(\mathring\omega)}$, provided
$0<\mathfrak Z=\int e^{-\mathscr S}Z_{X,\Theta(\mathring\omega)}d\nu^{\mathscr T}<\infty$
(\MT{prop:link-induced-action}). Finiteness of the generator prior alone does not imply
this second normalization. Gaussian determinants and their loop expansions belong to
\SG{ch:gauss-latent-links}.

\subsection{Recognition flow and learning the anchor}

\begin{keyequation}{First-order inference and the Gaussian Bayesian learning rule}
\[
\begin{aligned}
\dot Q&=-\operatorname{grad}^{\mathsf G}F,
&\frac{dF}{dt}&=-\lVert\operatorname{grad}^{\mathsf G}F\rVert_{\mathsf G}^2,\\
\mathring\Lambda'&=(1-\Delta t)\mathring\Lambda
+\Delta t\operatorname{pr}_{\mathring\Xi}\E_R\nabla^2\mathfrak i_\lambda,\\
\mathring\mu'&=\mathring\mu-\Delta t\mathring\Lambda'^{-1}\E_R\nabla\mathfrak i_\lambda.
\end{aligned}
\]
\end{keyequation}

Here $\mathsf G$ is the product of the Gaussian Fisher metric and the agents' recognition
metric, $\mathfrak i_\lambda$ is the held-agent free energy plus loop action and Gaussian
anchor negative log density, and $\operatorname{pr}_{\mathring\Xi}$ enforces the selected
full, edge-block, or diagonal precision family. The last two lines are explicit Euler in
Gaussian natural parameters, rewritten as mean and precision (\MT{prop:link-recognition-flow}).
They are the Gaussian Bayesian learning rule, standard in substance. The step is admissible
only when $\mathring\Lambda'\succ0$; finite-step descent also requires the stated relative
curvature bound in expectation coordinates and a sufficiently small step. For a positive
quadratic integrand in the full Gaussian family, a unit step gives its Gaussian minimizer.
Continuous descent has no nonconstant periodic orbit, and its nondegenerate critical points
have real linearized modes; convergence is not established in general.

Learning an anchor is a separate M-coordinate, changing a generative parameter at fixed
recognition. Its exact update matches the prior mean to $\E_R\mathring\omega$, not generally
the anchor itself to that mean (\MT{prop:link-anchor-m-coordinate}). At fixed covariance
this optimizer exists uniquely and includes the derivative of $\log Z_{\mathscr S}$.
With zero action the update is $\mathring\alpha^{\mathrm{new}}=\mathring\mu$; for a quadratic
loop action it is
$\mathring\alpha^{\mathrm{new}}=(I+\mathring\Sigma\sum_p\mathrm c_p\mathbf C_p^\top\mathbf C_p)
\mathring\mu$. Joint optimization of anchors and prior covariances can fail to attain a
maximum when loop terms are present. Natural-gradient recognition also supplies no automatic
belief--model--link rate separation: \MT{rem:link-rates} gives a held Gaussian sheet on which
all means relax at the same rate.

\subsection{Episode histories and Euclidean inertia}

A history prior multiplies the Gaussian base laws across $n=0,\ldots,n_1$ by
$\exp[-\sum_n\mathscr S(\mathring\omega^{(n)})
-\mathbf m\sum_n\mathfrak U(\mathring\omega^{(n+1)},\mathring\omega^{(n)})]$.
It is normalized for either nonnegative increment term: the Gaussian squared difference,
or the time-like loop defect of
$\exp(-\mathring\omega^{(n+1)}_e)\exp(\mathring\omega^{(n)}_e)$
(\MT{prop:link-history-prior}). Frames are static; the time-like links are identities,
not new latent variables. The equivalent forward kernels include backward normalizers
depending on the chosen horizon. Independently normalizing each increment generally gives
a different history law.

\begin{keyequation}{Inertia of the most probable generator history}
\[
\mathbf m\bigl(2\mathring\omega^{(n)}-\mathring\omega^{(n+1)}
-\mathring\omega^{(n-1)}\bigr)
+\nabla\mathfrak V_n(\mathring\omega^{(n)})=0,
\qquad 0<n<n_1.
\]
\end{keyequation}

For Gaussian increments this equation is exact at a critical history; $\mathfrak V_n$
contains the anchor, loop action, and any declared episode matter term
(\MT{prop:link-map-inertia}). Free ends have one-sided equations; held ends impose boundary
values. The sign corresponds to $\mathbf m\ddot{\mathring\omega}=+\nabla\mathfrak V$,
Newton's equation in potential $-\mathfrak V$, not real-time evolution in $+\mathfrak V$.
The increment precision is its mass. The time-like defect gives a generator-dependent
mass tensor, reducing to this scalar-mass quadratic at zero, with a bounded near-zero
correction to the equation. A one-dimensional most-probable path between the specified
wells is a monotone kink; that special-instance result does not establish an oscillation
or kink classification on nonabelian networks.

\subsection{The Proca-type operator and its approximation}

\begin{keyequation}{The Proca-type link equation: weak-current, near-zero form}
\[
\begin{aligned}
\mathbf K\mathring\omega&=\mathbf j,\\
\mathbf K&=\mathring\Sigma^{-1}
+\sum_p\mathrm c_p\mathbf C_p^\top\mathbf C_p
+\mathbf m(-\Delta_n)+\mathring\Phi,\\
\mathbf K&\succeq\mathrm c_{\mathring\Sigma}I,\qquad
\mathring\omega_{\mathrm{lin}}=\mathbf K^{-1}\mathbf j.
\end{aligned}
\]
\end{keyequation}

This is \MTeq{eq:link-linearized-equation}, with the operator of
\MT{prop:link-linearized-equation}. The anchor precision supplies a bare mass; the cycle
sum supplies curl-curl on the agent multigraph; $-\Delta_n$ is the positive episode
Laplacian, including the free-end terms; $\mathring\Phi$ is the positive matter quadratic
form built from $W_a$ and the source second moments. Its source is the sum of projected
zero-generator currents, weighted by the fixed $\beta_a$. The smallest eigenvalue
$\mathrm c_{\mathring\Sigma}>0$ bounds the inverse and gives a unique linear solution
on every finite graph, not only a small network.

The exact Hessian at zero is
\[
\nabla^2\mathscr S^{\mathrm{hist}}(0)=\mathbf K+\mathbf N_0,
\qquad
\langle\mathring\omega,\mathbf N_0\mathring\omega\rangle
=-\sum_{n,a}\beta_a\Tr\bigl(\mathbf J_a^\top
(\mathrm d\rho\mathring\omega^{(n)}_e)^2\bigr).
\]
Thus $\mathbf K\mathring\omega=\mathbf j$ keeps first order jointly in generators and
currents. At fixed current, the generator-linear Taylor equation includes $\mathbf N_0$;
cancellation of the total source does not by itself make the individual-current correction
small. The nonlinear comparison is proved with product Gaussian matter, fixed
generator-independent weights, zero anchor, and the specified quadratic comparison records.
If $\mathrm c_0=\lVert\mathbf N_0\rVert$ and
$\mathrm c_{\mathrm r}=\sup_{\lVert\mathring\omega\rVert\leq\mathrm r}
\lVert D^3\mathscr S^{\mathrm{hist}}\rVert$ satisfy
$\mathrm c_0+\mathrm c_{\mathrm r}\mathrm r\leq\mathrm c_{\mathring\Sigma}/2$ and
$\lVert\mathbf j\rVert\leq\mathrm c_{\mathring\Sigma}\mathrm r/4$, there is exactly one
critical history in that ball and its discrepancy from the linear solution obeys
\[
\lVert\mathring\omega_\star-\mathring\omega_{\mathrm{lin}}\rVert
\leq\frac{2\lVert\mathbf j\rVert}{\mathrm c_{\mathring\Sigma}^2}
\left(\mathrm c_0+\frac{\mathrm c_{\mathrm r}\lVert\mathbf j\rVert}
{\mathrm c_{\mathring\Sigma}}\right).
\]
At fixed third-derivative bound this is second order in the individual currents. It is
a local conclusion, not uniqueness of all stationary histories.

\begin{analogy}{structural}
The boxed operator has the mass, curl-curl, and positive temporal-difference form of a
massive Euclidean vector operator of Proca type. Its components live in $\mathfrak g$ on
agent edges; its ``temporal'' direction is the inference episode index $n$. No continuum
limit, physical propagation law, or Lorentzian field equation is asserted. Nonabelian
brackets enter the higher-order terms, while the graph curl-curl and episode terms act
componentwise at the displayed order. History inertia and first-order recognition
relaxation remain different statements.
\end{analogy}

\section{What is proved, and how solid it is}
\label{comp:links-status}

\begin{scorecard}
\MT{def:link-generator-prior} & Anchors, covariances, metric, and penalized walks specify the latent-link model. & \CHOICE & They are supplied generative data, not consequences of the bound.\\
\MT{prop:link-generator-gauge} & Passive frames leave generators fixed; reference changes act by one adjoint map. & \SOLID & Exact coordinate laws.\\
\MT{prop:link-chart-image} & Exact exponential-image restriction; suitable loop products reach $\GL^+(K)$. & \SOLID & Excludes nonexponential edge twists.\\
\MT{cor:link-no-invariant-defect} & Continuous class functions cannot separate unipotent holonomy from identity. & \SOLID & General $\GL(K,\R)$ and $\GL^+(K)$, $K\geq2$; a negative result.\\
\MT{prop:link-generator-prior}, \MT{prop:link-loop-defect} & Prior normalization, metric loop measurement, and covariance under carried reference data. & \SOLID & Compact or noncompact $G$; fixed-matrix symmetry is narrower.\\
\MT{prop:link-latent-normalization}, \MT{prop:link-gaussian-recognition-vfe} & Unit-mass joint, conditional decomposition, and finite Gaussian link sector. & \SOLID & Any finite record hypergraph; decomposition assumes bounded records and positive conditional evidence.\\
\MT{prop:link-field-equation}, \MT{cor:link-delta-rule} & Gaussian mean and precision stationarity; adjoint-filtered expected delta current. & \SOLID & Product generator--matter recognition, dominated derivatives; stationarity, not uniqueness.\\
\MT{cor:link-prior-necessary} & Singular product-matter regression optimum and stabilizer degeneracy. & \SOLID & Chart critical points and correlated matter are explicit exceptions to broader slogans.\\
\MT{thm:link-held-limit} & Narrowing the prior recovers held-link infima, concentration, and conditionally evidence. & \COND & Bounded records, positive conditional evidence, lower semicontinuity, and an approximating class; evidence also needs continuity.\\
\MT{prop:link-conditional-transfer} & Held-link conclusions transfer conditionally; mixtures need not preserve them. & \SOLID & The property must be mixture-stable for a marginal conclusion.\\
\MT{prop:link-induced-action} & Jointly normalized matter induces $-\log Z$; record extension induces none. & \COND & Joint undirected normalizer must be finite; generator normalization does not ensure it.\\
\MT{prop:link-recognition-flow} & Fisher descent and Bayesian learning rule as its Euler step. & \SOLID & Standard in substance; discrete admissibility and descent have separate hypotheses.\\
\MT{prop:link-anchor-m-coordinate} & Unique exact anchor update at fixed covariance; joint covariance update can fail. & \SOLID & Prior normalizer included; no unconditional joint M-step existence.\\
\MT{prop:link-history-prior} & Normalized Gibbs history and exact horizon-dependent forward kernels. & \SOLID & Finite horizon, static frames, no time-like variables.\\
\MT{prop:link-linearized-equation} & Positive Proca-type operator and a controlled nonlinear local solution. & \COND & Small individual currents and near-zero ball; fixed product Gaussian matter and declared records.\\
\end{scorecard}

The attainment clauses of \MT{cor:link-prior-necessary} are \COND: they need the stated
lower-bounded matter objective or the specified nonnegative-record Gaussian product model.
\MT{prop:link-map-inertia} is \SOLID\ for its general stationarity equations and coercive
existence statement, with coercivity required for existence; its scalar monotone-kink clause
is \TOY. The main prior, coordinate laws, field equation, and history construction apply
to arbitrary finite agent counts. Translation instances and the scalar kink ground those
constructions; they do not set the direction for large nonabelian networks. None of these
results supplies a thermodynamic limit or a universal dynamical phase classification.

\begin{hygiene}
Substantive commitments are the exponential-image restriction; the supplied fiber metric,
anchors, covariances, and loops; product generator--matter recognition for the displayed
Gaussian field equations; mixture stability for marginal transfers; finiteness of the
joint undirected normalizer; static frames and a chosen finite episode horizon; lower
semicontinuity and the approximating class in the held limit; and the small-current,
near-zero hypotheses of the nonlinear Proca comparison. A proposed finite Euler step must
preserve positive precision and meet its relative curvature bound to inherit descent.
These conditions cannot be replaced by the fact that a generator prior is normalized.

Regularity hygiene includes standard Borel spaces for conditional laws, measurable factor
kernels, finite Gaussian moments, and differentiation under integrable growth bounds.
The smooth finite-dimensional matrix-exponential comparison models satisfy the stated
growth mechanism. A different matter model still has to meet those bounds. Normalization
of the prior is unconditional within its declared construction; convergence of recognition
is a separate question.
\end{hygiene}

\section{Open ends}
\label{comp:links-open}

\MT{open:link-coarse-closure} asks for a finite-dimensional weakly continuous injective
parameter family of laws on $\GL^+(K)$ containing every exponential image of a nondegenerate
Gaussian and closed under convolution, or a proof that such a family cannot exist. Exact
coarse pushforwards exist without that closure. Products of exponential chart links may
have no single real logarithm, so writing fine links with generators does not automatically
give coarse links the same coordinates or law family.

The interaction graph remains structural; \MT{open:obs-edge-creation} concerns its creation
and deletion. A slow-manifold reduction with generator recognition means in the slow block
remains within \MT{open:obs-slow-manifold}. The chapter does not establish general
recognition convergence, uniqueness outside the controlled ball, prediction gains, or
existence of a joint anchor-and-covariance M-coordinate with loops. These are limits of
what is delivered, not consequences of the Proca analogy.

\begin{verdictbox}{\DELIVERED}
\textbf{Claimed.} The technical chapter turns independent links into latent exponential
generators with a normalized prior, includes them in the collective evidence bound,
derives their field equation and recognition flow, recovers held links in a narrowing
limit, and supplies a normalized episode history with inertia and a controlled
Proca-type linearized operator.

\textbf{Delivered.} These constructions and their stated equations are established at
one base point, including noncompact structure groups. The Gaussian recognition step
is the Bayesian learning rule in standard substance; the history inertia is an exact
stationarity equation for Gaussian increments; the positive linear operator has a
unique solution, and its nonlinear comparison is local under explicit small-current
conditions. The negative results about class functions and unregularized links explain
why the metric and prior are substantive parts of the model.

\textbf{Gap.} Closure of a finite-dimensional coarse link-law family remains open.
The chapter claims no physical Proca field, no general recognition convergence, no
global uniqueness of nonlinear histories, and no predictive advantage over held links.
Those limits are explicit in its scope. Calling links dynamic is justified for their
recognition and episode laws; physical inertia, spontaneous rate separation, and closed
coarse dynamics do not follow from those laws.
\end{verdictbox}
